\section{Squared deletion delays}
\label{sec:delay}

We now combine the root-test probability with the remaining-duration bound.
The two estimates use compatible filtrations: $A_m$ is known under
$\cH_m$, while the upper bound $W_m$ is already measurable under
$\cG_m$ and therefore does not reveal the hidden root tests.

\begin{proposition}[Conditional squared-delay bound]
\label{prop:conditional-delay}
For every $1\le i\le L$,
\begin{align}
    \E[D_i^2\mid\cE]
    \le{}&
    2K_0^2B_{\mathrm{t}}^d\,d^{2\eta}\ind_{\{V=0\}}
    \notag\\
    &+2K_0LB_{\mathrm{t}}^d\,d^\eta\ind_{\{V=1\}}
    +L^2\ind_{\{V\ge2\}}.
    \label{eq:conditional-delay-bound}
\end{align}
\end{proposition}

\begin{proof}
For $i=L$, the delay is zero.  Otherwise,
\eqref{eq:delay-count} implies
\[
    D_i^2
    \le
    2\sum_{m=i}^{L-1}\ind_{A_m}
       \sum_{s=m}^{L-1}\ind_{A_s}.
\]
Condition first on $\cH_m$, where $A_m$ is known, and apply
Corollary~\ref{cor:remaining-duration}.  Then condition on $\cG_m$, where
$W_m$ is known.  Conditioning the resulting estimates down to $\cE$ gives
\begin{equation}
    \E[D_i^2\mid\cE]
    \le
    2\sum_{m=i}^{L-1}
      \E\bigl[
        W_m\Prb(A_m\mid\cG_m)
        \mid\cE
      \bigr].
    \label{eq:tower-delay}
\end{equation}

On $\{V=0\}$, Lemma~\ref{lem:root-tests} and
\eqref{eq:W-definition} bound the product inside the expectation by
\[
    K_0B_{\mathrm{t}}^d\,d^\eta
    \ind_{\{S_m\ne\varnothing\}}p_m^{k-\eta}.
\]
For every $k\ge3$,
\begin{equation}
    k-\eta\ge\eta,
    \label{eq:critical-exponent}
\end{equation}
because $k\ge2k/(k-1)$.  Since $0\le p_m\le1$, the preceding display is at
most
\[
    K_0B_{\mathrm{t}}^d\,d^\eta
    \ind_{\{S_m\ne\varnothing\}}p_m^\eta.
\]
Lemma~\ref{lem:occupation}, applied from time $i$ and then conditioned down
to $\cE$, bounds the sum in \eqref{eq:tower-delay} by
$2K_0^2B_{\mathrm{t}}^d\,d^{2\eta}$.

On $\{V=1\}$, the corresponding product is at most
\[
    K_0B_{\mathrm{t}}^d\,d^\eta p_m^{k-1-\eta}
    \le K_0B_{\mathrm{t}}^d\,d^\eta,
\]
because $k-1-\eta\ge0$ for $k\ge3$.  There are at most $L$ terms.  If
$p_m=0$, the conditional pivotality probability is zero by
Lemma~\ref{lem:root-tests}, so no division by zero occurs.  Finally, on
$\{V\ge2\}$ use $D_i\le L$.
\end{proof}

It remains to average the factors in
\eqref{eq:conditional-delay-bound}.  The clause coordinates need not have
the same law, but the local incidence hypotheses give uniform incidence
probabilities.

\begin{lemma}[Incidence and collision moments]
\label{lem:incidence-moments}
Uniformly in $n\ge2k$, $i\le L$, and the conditioned value of $C_i$,
\begin{align}
    \E[B_{\mathrm{t}}^d\,d^{2\eta}]&=O(1),
    \label{eq:exp-moment}\\
    \E[B_{\mathrm{t}}^d\,d^\eta\ind_{\{V\ge1\}}]
      &=O(n^{-1}),
    \label{eq:one-collision-moment}\\
    \Prb(V\ge2)&=O(n^{-2}).
    \label{eq:two-collision-prob}
\end{align}
The implicit constants depend only on
$k,\rho,B_{\mathrm{t}},\gamma_1,\gamma_2$.
\end{lemma}

\begin{proof}
For $j\ne i$, let $I_j$ be the indicator that $C_j$ meets at least one
root, and let $J_j$ be the indicator that it meets at least two roots.
The local incidence assumptions give
\begin{equation}
    \Prb(I_j=1)\le \frac{\gamma_1}{n},
    \qquad
    \Prb(J_j=1)\le \frac{\gamma_2}{n(n-1)}.
    \label{eq:incidence-probabilities}
\end{equation}
The indicators at distinct indices are independent, although their
probabilities may differ.  Consequently, for every fixed $t\ge0$,
\[
    \E e^{tZ}
    =\prod_{\substack{j\le L\\j\ne i}}
       \left(1+(e^t-1)\Prb(I_j=1)\right)
    \le
    \exp\bigl((e^t-1)\gamma_1\rho\bigr).
\]
Choosing $t>\log B_{\mathrm{t}}$ and absorbing the polynomial factor
$d^{2\eta}$ into a slightly larger exponential proves
\eqref{eq:exp-moment}.

Conditional on $J_j=1$, the other clause coordinates retain their original
independent laws, and $Z$ is one plus the incidence count from those other
coordinates.  The same exponential-moment estimate therefore gives
\[
    \E[B_{\mathrm{t}}^d\,d^\eta\mid J_j=1]=O(1)
\]
uniformly in $j$.  Since
$\ind_{\{V\ge1\}}\le\sum_{j\ne i}J_j$,
\[
    \E[B_{\mathrm{t}}^d\,d^\eta\ind_{\{V\ge1\}}]
    \le
    O(1)\sum_{j\ne i}\Prb(J_j=1)
    =O(n^{-1}),
\]
which proves \eqref{eq:one-collision-moment}.  Finally, independence
across indices gives
\[
    \Prb(V\ge2)
    \le \E\binom V2
    =\sum_{j<\ell}\Prb(J_j=1)\Prb(J_\ell=1)
    =O(n^{-2}).
\]
\end{proof}

Under the pointwise comparison hypothesis, the concrete bounds
\eqref{eq:local-parameters} follow because, under $U_{n,k}$, a specified
root is included with probability $k/n$ and a specified root pair with
probability $k(k-1)/[n(n-1)]$; union bounds over roots and root pairs and
the upper density factor $b$ give \eqref{eq:incidence-probabilities}.

\begin{proposition}[Uniform deletion-delay moment]
\label{prop:uniform-delay}
For every fixed $\rho<\infty$,
\[
    \sup_{\substack{n\ge2k,\,1\le L\le\rho n\\1\le i\le L}}
    \E D_i^2
    <\infty,
\]
where the supremum is uniform over all clause-law arrays satisfying the
local hypotheses of Theorem~\ref{thm:local-criterion}.
\end{proposition}

\begin{proof}
Average \eqref{eq:conditional-delay-bound}.  The first term is $O(1)$ by
\eqref{eq:exp-moment}.  The second is
\[
    O(L)\,
    \E[B_{\mathrm{t}}^d\,d^\eta\ind_{\{V\ge1\}}]
    =O(Ln^{-1})=O(1),
\]
and the third is
\[
    L^2\Prb(V\ge2)=O(L^2n^{-2})=O(1).
\]
The bounds were uniform in the conditioned value of $C_i$, so that
conditioning may be removed.
\end{proof}

Combining Proposition~\ref{prop:uniform-delay} with
Lemma~\ref{lem:coordinate-omission} proves
\begin{equation}
    \Var(H_n\wedge L)=O(n)
    \qquad(n\ge2k,\ L\le\rho n),
    \label{eq:capped-variance}
\end{equation}
with the dependence asserted in Theorem~\ref{thm:local-criterion}.  For the
finitely many sizes $k\le n<2k$, the deterministic bound
$H_n\wedge L\le L\le\rho n$ gives the same $O(n)$ conclusion after
increasing the constant.  This proves the clipped part of the local
criterion.
